\section{Introduction}\label{sec:intro}

Model checking a real-time system is expensive. A network of timed automata denotes an
uncountable state space, and the standard remedy --- collapsing clock valuations into convex
\emph{zones} --- leaves a symbolic state space that still grows exponentially in the number of
components. Industrial-scale verification therefore routinely runs for hours or exhausts
memory, and since single-thread performance stopped improving the only remaining source of
capacity is parallelism. The question is not whether to parallelise timed reachability, but
what exactly to parallelise.

The obvious answer is wrong. Successor generation is embarrassingly parallel, so one might
expect a straightforward speed-up from expanding many states at once. In practice the
bottleneck is not generating successors but \emph{storing} them. An untimed model checker
keeps a hash set of visited states, and a lock-free hash table settles the question. A timed
model checker cannot: two symbolic states sharing a discrete configuration must be compared
by zone inclusion, either may absorb the other, and a newly discovered zone can render several
stored zones redundant at once. The visited set is thus a family of per-configuration
\emph{antichains} that grow and shrink throughout the search. Every design decision in a
parallel timed model checker is, in the end, a decision about how threads share that
structure, and it is there that existing approaches run into trouble.

Prior multi-core work demonstrates that the idea is sound while exposing three specific
difficulties. First, \emph{memory}. Sharing symbolic states across threads normally means
materialising a zone matrix per stored state, and at the scale where parallelism is worth
having this dominates the footprint; we measure a reference multi-core implementation using
$13.5$\,GB where ours uses $1.0$\,GB on the same model. Second, \emph{inconsistency across
benchmarks}. Reported speed-ups are excellent on some models and collapse on others, and the
published explanation --- that scalability degrades as the ratio of symbolic to discrete
states rises --- identifies the symptom without removing it. Third, and most awkwardly,
\emph{parallel is not always better}: on models with small discrete parts a parallel checker
can lose outright to sequential \textsc{Uppaal}, because the parallel implementation explores
a larger state space than the sequential tool it is compared against. A tool that is
sometimes slower than the baseline it aims to replace is difficult to deploy.

This paper presents a parallel reachability algorithm addressing all three. We share the
visited set across threads but make sharing cheap enough that expansion and insertion can be
\emph{fused}: the thread that generates a successor stores it immediately, with no
intermediate buffer and no routing phase. We keep memory low by storing each distinct zone
once and referring to it by index --- justified by the observation that a run over $2.7$
million symbolic states may contain only $120$ distinct zones --- and by packing the discrete
configuration into a key from which it can be recovered, so it costs nothing per stored state.
We then show that the fused design has a precise weakness: its critical section is
proportional to antichain length, hence to the very ratio prior work identified, and we
remove it with a second engine that defers subsumption to a lock-free per-level phase. Because
we never trade state-space size for speed, our checker explores exactly the state space
\textsc{Uppaal} explores on the models where the two are comparable, so a speed-up is a
speed-up rather than an artefact of doing different work.

Concretely, we contribute: a fused expansion/insertion engine over a shared lock-striped
visited set, together with the interning and key-packing that make it affordable
(\cref{sec:approach}); an analysis of its critical section and a batched engine that reduces
it from $\Theta(\rho)$ to $O(1)$, where $\rho$ is the mean antichain length
(\cref{sec:batching}); and an evaluation against \textsc{Uppaal} and \textsc{opaal}+\ltsmin{}
under a uniform time and memory budget (\cref{sec:eval}), in which we are faster than both
wherever both complete, use one to two orders of magnitude less memory, and reproduce
\textsc{Uppaal}'s state counts exactly.
